\documentclass[11pt,a4paper]{article}

\usepackage[margin=2.3cm]{geometry}
\usepackage[T1]{fontenc}
\usepackage[utf8]{inputenc}
\usepackage{lmodern}
\usepackage{amsmath,amssymb,bm}
\usepackage{booktabs}
\usepackage{enumitem}
\usepackage{xcolor}
\usepackage{listings}
\usepackage{hyperref}

\hypersetup{
  colorlinks=true,
  linkcolor=blue,
  urlcolor=blue
}

\definecolor{codegray}{RGB}{245,245,245}
\definecolor{codeblue}{RGB}{0,70,140}

\lstset{
  language=Python,
  basicstyle=\ttfamily\small,
  backgroundcolor=\color{codegray},
  keywordstyle=\color{codeblue},
  frame=single,
  breaklines=true,
  showstringspaces=false
}

\title{\textbf{Practical Session: Locomotion Software for Legged Robots, real experiments}}
\author{}
\date{}


\begin{document}

\maketitle


\section*{Objectives}
This practical session introduces the main components of a model-based locomotion software stack for a quadruped robot in simulation. 
There are several blocks/modules  in a locomotion framework, we will do exercises separately on  them and then put all the modules together and experience and do tests with the whole framework.
The goal is to progressively build, tune, and test the following modules of the framework: Swing controller, Stance controller, Foothold generator, Gait scheduler. 


\section*{Install the software}
\textbf{First repo}: 

\begin{verbatim}
git stash
git fetch origin
git checkout origin/feat/lesson_saturday
\end{verbatim}

\noindent Very important: We now use the branch \texttt{feat/lesson\_saturday}.\\
Very important: Every time you open a new terminal, you need to activate the conda environment first, see the README!



\section{Exercise 1: Swing-Foot Trajectory Generator and Control}

The swing-foot trajectory must move the foot smoothly from lift-off to the desired touchdown point. It should provide sufficient clearance above the terrain and have low touchdown velocity.

For a cubic trajectory:

\begin{align*}
p(t) = a_0 + a_1t + a_2t^2 + a_3t^3.
\end{align*}

Choose the coefficients to satisfy initial and final position and velocity boundary conditions.
The file to modify is \verb|Quadruped-PyMPC/quadruped_pympc/helpers/cubic_swing_trajectory.py| 
Startin and final position of the spline are already given in the main of this file, that you can launch with
 \verb|python3 Quadruped-PyMPC/quadruped_pympc/helpers/cubic_swing_trajectory.py| 

Hint: search for EXERCISE 1.1-1.4 in the file and implement between the dotted lines.

\begin{enumerate}[label=\arabic*.]
    \item Generate horizontal motion from lift-off (initial point) to touchdown (final point).
    \item Add a vertical swing-height profile - the robot should not slide the feet on the terrain!
    \item Compare position, velocity, and acceleration continuity.
    \item Verify that touchdown velocity is close to zero. Real robot can suffer from discontinuities!
\end{enumerate}



For tracking the trajectory, a control law is needed (1):

\begin{align*}
\bm{\tau}_{sw,i}  =
\bm{J}_i^T \left[ \bm{K}_p
\left(
\bm{p}^{\mathrm{ref}}-\bm{p}
\right)
+
\bm{K}_d
\left(
\dot{\bm{p}}^{\mathrm{ref}}-\dot{\bm{p}}
\right)\right].
\end{align*}
 

A better control law can be obtained by compensating the gravitational effects and coriolis (2).  

\begin{align*}
\bm{\tau}_{sw,i}  =
\bm{J}_i^T \left[ \bm{K}_p
\left(
\bm{p}^{\mathrm{ref}}-\bm{p}
\right)
+
\bm{K}_d
\left(
\dot{\bm{p}}^{\mathrm{ref}}-\dot{\bm{p}}
\right)\right] + h(\bm{q}, \dot{\bm{q}}).
\end{align*}


The file to modify is \verb|Quadruped-PyMPC/quadruped_pympc/helpers/swing_trajectory_controller.py|.
Run its main with \verb|python3 Quadruped-PyMPC/quadruped_pympc/helpers/swing_trajectory_controller.py| 
HINT: search for EXERCISE 1.5-1.7 in the file and implement between the dotted lines.
\begin{enumerate}[label=\arabic*.]
    \item Implement controller (1), and check the tracking of the spline trajectory; Tune the gains \(K_p\) and \(K_d\) to improve tracking - range [200-500 Kp], [5-20 Kd].
    \item Implement controller (2), and check the tracking of the spline trajectory; Tune the gains \(K_p\) and \(K_d\) to improve tracking - range [200-500 Kp], [5-20 Kd].
    \item These are the controllers you will use on the real robot, try to obtain a smooth action!
\end{enumerate}
KP and KD can be tuned in The file to modify is \verb|Quadruped-PyMPC/quadruped_pympc/config.py|, lines 202-203.


\section{Exercise 2: Stance Controller}

During stance, the desired contact force computed by the MPC is converted into joint torques through the foot Jacobian:

\begin{align*}
\bm{\tau}_{st,i} = -\bm{J}_i^T \bm{f}_i^{\mathrm{des}}.
\end{align*}

The Jacobian must be consistent with the frame in which the contact force is expressed.

Modify \verb|Quadruped-PyMPC/quadruped_pympc/interfaces/wb_interface.py|, lines 372-377, search for "EXERCISE 2.1"
\begin{enumerate}[label=\arabic*.]
    \item Extract the foot Jacobian from the robot model.
    \item Take the GRF from the MPC.
    \item Compute the joint torques using the Jacobian and the desired contact force.
\end{enumerate}

% \section{Raibert-Style Foothold Generator}
% 
% At lift-off, the desired touchdown position can be computed from the nominal hip location, desired velocity, and velocity tracking error:
% 
% \begin{align*}
% \bm{p}_{f,i}^{\mathrm{des}} =
% \bm{p}_{\mathrm{hip},i}
% +
% \bm{p}_{i}^{\mathrm{nom}}
% +
% \frac{T_{\mathrm{st}}}{2}\bm{v}^{\mathrm{des}}
% +
% k_v
% \left(
% \bm{v} - \bm{v}^{\mathrm{des}}
% \right),
% \end{align*}
% 
% where \(T_{\mathrm{st}}\) is the stance duration and \(k_v\) is the velocity-feedback gain.
% 
% 
% \subsection*{Exercises}
% 
% 
% 
% \begin{enumerate}[label=\arabic*.]
%     \item Implement the nominal foothold generator.
%     \item Command forward, lateral, and turning motion.
%     \item Visualize the planned touchdown locations.
%     \item Tune \(k_v\) and observe its influence on velocity recovery.
%     \item Add reachable-workspace bounds around the hip.
% \end{enumerate}




\section*{Exercise 3: Gait Scheduler}

The gait scheduler generates a binary contact state for each foot:

\begin{align*}
	s_i(t) \in \{0,1\},
\end{align*}

where \(s_i=1\) denotes stance and \(s_i=0\) denotes swing.

Each gait is described by its cycle duration, duty factor, and phase offsets.

\begin{table}[h]
	\centering
	\begin{tabular}{lll}
		\toprule
		Gait & Swing pattern & Typical use \\
		\midrule
		Crawl & One foot swings at a time & Slow and statically stable walking \\
		Pace & Left and right lateral pairs alternate & Fast lateral-pair gait \\
		Trot & Diagonal pairs alternate & Dynamic walking and running \\
		Bound & Front and hind pairs alternate & Fast locomotion \\
		\bottomrule
	\end{tabular}
\end{table}


The file to modify is \verb|Quadruped-PyMPC/quadruped_pympc/helpers/periodic_gait_generator.py|
You can launch its main with \verb|python3 Quadruped-PyMPC/quadruped_pympc/helpers/periodic_gait_generator.py|
You can select the gait in \verb|Quadruped-PyMPC/quadruped_pympc/config.py| in 217. Remember to change duty factor and step frequency in the config file as well there.
\begin{enumerate}[label=\arabic*.]
	\item Generate the contact schedule for the crawl.
	\item How many different types of crawl can you generate?
	\item Change duty factor and gait frequency during simulation to find a stable gait.
\end{enumerate}



\section*{Exercise 4: Real robot}

With all you changes done in these exercises, try to run the simulation:
\verb|python3 Quadruped-PyMPC/simulation/simulation.py|

If it works, congrats, we can try it on the real robot!



% \section{Integral Action and Payload Adaptation}
% 
% Integral action can compensate for persistent tracking errors caused by unmodelled disturbances, such as an additional payload:
% 
% \begin{align*}
% \bm{e}_I(t)
% =
% \int_{0}^{t}
% \left(
% \bm{v}^{\mathrm{des}}-\bm{v}
% \right)
% dt.
% \end{align*}
% 
% The integral state can be added to the body-position or velocity controller.
% 
% \subsection*{Experiment}
% 
% \begin{enumerate}[label=\arabic*.]
%     \item Tune the controller for the nominal robot.
%     \item Add an additional 5kg mass to the robot model in MuJoCo.
%     \item Test walking without integral action.
%     \item Activate the integral term and compare velocity tracking.
% \end{enumerate}

%
%\subsection*{Final Experiment}
%
%\begin{enumerate}[label=\arabic*.]
%    \item Create a rough-terrain scenario in MuJoCo.
%    \item Run the robot with terrain estimation disabled.
%    \item Repeat with terrain estimation enabled.
%    \item Compare base orientation, foot tracking error, contact forces, and velocity tracking.
%    \item Discuss the role of terrain estimation in stability and locomotion robustness.
%\end{enumerate}

%\section{Expected Final Demonstration}
%
%At the end of the practical session, the robot should be able to:
%
%\begin{itemize}
%    \item Walk using crawl, pace, and trot gaits;
%    \item Track a commanded body velocity using centroidal MPC;
%    \item Generate swing-foot trajectories and stance torques;
%    \item Use numerical or QP-based inverse kinematics;
%    \item Adapt footholds with Raibert feedback and foothold optimization;
%    \item Recover from pushes through corrective stepping;
%    \item Compensate partially for an added payload using integral action;
%    \item Walk on uneven terrain using terrain-aware foothold and base-orientation references.
%\end{itemize}

\end{document}
